\documentclass[10pt,letterpaper]{article}
\usepackage{cornell-notes}

\title{Clause 14.06.02: The Std Terminate Function}
\author{}
\date{}
\setDocTitle{Clause 14.06.02: The Std Terminate Function}
\setDocSubtitle{C++ 2024}
\setDocOwner{C++ 2024 Cornell Notes}
\setDocDate{\today}
\setCornellCollection{C++ 2024 Cornell Notes}
\setCornellUnitType{Clause}
\setCornellUnitNumber{14.06.02}
\setCornellUnitTitle{The Std Terminate Function}
\hypersetup{pdftitle={Clause 14.06.02: The Std Terminate Function},pdfsubject={Cornell notes},pdfauthor={}}

\begin{document}
\maketitle
\begin{CornellOverviewBox}{Overview}
\textbf{Classification:} Core language - Normative\\[2pt]
\textbf{Learning objective:} Understand the rules, terminology, and semantic consequences associated with The std::terminate function.
\end{CornellOverviewBox}\begin{CornellSummaryBox}{Summary}
{\sffamily\bfseries\color{Accent} Summary}\par\vspace{3pt}Section 14.6.2 focuses on The std::terminate function. The section supplies the rules and semantic distinctions needed to interpret this topic in context with its cross-references. When applying it, preserve the distinction between mandatory requirements, recommendations, permissions, and behavior deliberately left undefined, unspecified, or implementation-defined.
\end{CornellSummaryBox}

\end{document}
